\chapter{Att kontrollera ett påstående}
\label{app:verifiera}

Det här materialet gör, som allt undervisningsmaterial, en rad
påståenden.  De flesta går att pröva vid datorn: att en modul körs när
den importeras ser ni själva genom att köra \texttt{use\_hello.py}, och
att en fil som heter \texttt{statistics.py} skuggar modulen med samma
namn ser ni genom att döpa om en fil.  Men några av påståendena handlar
om vad andra har sagt eller visat, och dem kan ingen körning avgöra:
\begin{enumerate}
  \item att det är bra att inte upprepa sig i ett program --- principen
    \emph{DRY} --- och att den kommer från Hunt och Thomas
    (\cref{sec:tva-program});
  \item att en modul, liksom en funktion, bör ha ett enda ansvar ---
    \emph{SRP} --- och att principen är Robert C. Martins
    (\cref{sec:modul});
  \item att den enklaste lösning som duger är att föredra ---
    \emph{KISS} --- och att uttrycket tillskrivs Kelly Johnson
    (\cref{sec:kiss});
  \item att Python gör det materialet säger att det gör: att en modul
    är en fil, att \mintinline{python}{import} kör den, att programmets
    katalog söks före standardbiblioteket, och så vidare
    (\cref{sec:modul,sec:main,sec:stdlib,sec:pypi});
  \item att studenter bygger egna regler för var ett namn letas upp,
    vilket \cref{sec:modul} är utformat mot.
\end{enumerate}
Hur vet man om sådant stämmer?  Samma fråga möter er varje gång ni läser
en lärobok, en webbsida eller ett svar från en språkmodell --- och varje
gång ni själva ska skriva en rapport.  Den här bilagan sammanfattar
arbetsgången och visar i \cref{tab:paastaenden} var vart och ett av
påståendena är kontrollerat.  Arbetsgången beskrivs utförligare i
bilagan med samma namn i föreläsningen \emph{Algoritmiskt tänkande}, och
de tre principerna har redan prövats mot litteraturen i tidigare
föreläsningar; det görs inte om här.

\section{Gör påståendet till en fråga}

Ett påstående går att kontrollera först när det är tydligt vad som skulle
kunna visa att det är \emph{fel}.  Formulera det därför som en fråga med
ett svar.  \enquote{Kommer DRY från Hunt och Thomas?} kan besvaras med ja
eller nej; \enquote{moduler är viktiga} kan det inte, eftersom
\enquote{viktigt} är ett omdöme och inte ett faktum.  Ofta döljer sig
flera påståenden i ett: \enquote{SRP är Martins etablerade princip} är
två frågor --- \emph{vems} är den, och är den \emph{etablerad} --- och
de kräver olika slags källor.

\section{Välj rätt sorts källa}

Olika frågor besvaras av olika källor.
\begin{description}
  \item[Uppslagsverk] för vad ett ord \emph{betyder}.
  \item[Primärkällan] för var en idé \emph{kommer ifrån}, och för ett
    specifikt resultat.  För vad ett programspråk \emph{gör} är
    primärkällan språkets egen dokumentation --- och programmet självt:
    ett påstående om vad Python gör kontrolleras genom att köra Python.
  \item[Översikter] (\foreignlanguage{english}{reviews}) för om något är
    \emph{etablerat}, eller för vad forskningen \emph{sammantaget}
    säger.
  \item[Enskilda studier] för ett specifikt resultat --- med förbehållet
    att en enskild studie bara visar att något observerats en gång, i en
    viss population.
\end{description}
Var man hittar dem: universitetsbibliotekets söktjänst, och databaser som
Scopus, Web of Science, IEEE Xplore, DBLP (datalogi) och OpenAlex
(öppen).  I kursens bilagor används kommandoradsverktyget
\texttt{scholar}\footnote{Paketet \texttt{scholarcli} på PyPI:
  \url{https://pypi.org/project/scholarcli/}.}, som ställer samma fråga
till flera databaser på en gång och sparar allt som hände i en
\emph{session}.  Det är för övrigt ett paket av just det slag
\cref{sec:pypi} handlar om.

\section{Sök så att du kan ha fel}

Den som söker på namn hon redan känner till hittar det hon väntar sig.
Sök därför \emph{ämnesdrivet} --- på begrepp, inte författare.  Några
praktiska regler:
\begin{itemize}
  \item Håll frågorna korta.  Flera korta frågor slår en lång.
  \item Ställ \emph{samma} frågor i alla databaser, översatta till varje
    databas frågespråk.
  \item Sök också efter \emph{invändningar}: lägg till ord som
    \enquote{critique}, \enquote{limitations} eller \enquote{justified}.
    Ett påstående som överlevt en motsökning är värt mer än ett som
    aldrig prövats, och en invändning som hittas blir ett
    \emph{förbehåll} i svaret, inte ett nederlag.  Bilaga C i
    föreläsningen \emph{Funktioner} är ett sådant fall: motsökningen
    visade att KISS fanns i tryck före den som brukar få äran, och
    därför står det \emph{tillskrivs} i \cref{sec:kiss}.
  \item Räkna med att databasernas rankning missar somligt.  Då hämtas
    källan som \emph{känt dokument}: man vet att den finns och slår upp
    den direkt.  Det är i sin ordning, bara det sägs --- alla
    Pythonkällor i det här materialet är hämtade så.
\end{itemize}

\section{Läs i fulltext, och skriv ner hur}

En träff i en databas är inte ett belägg.  Sammanfattningen
(\foreignlanguage{english}{abstract}) räcker inte heller: den säger vad
författarna vill ha sagt, inte vad de visat.  Läs hela texten och leta
upp \emph{det ställe} som styrker påståendet --- ett ordagrant citat med
sidnummer.  Hittar ni inget sådant ställe stöder källan inte påståendet,
hur passande titeln än är.  Räkna aldrig en källa ni inte läst.  Det
gäller också källor ni ärvt: en post som redan står i litteraturlistan
är inte kontrollerad förrän ni själva läst stället.  Fleurys studie av
vilka regler nybörjare bygger för var ett namn letas upp fanns i kursens
sammanställning med bara sammanfattningen läst; för det här materialet
lästes hela artikeln, och citatet i \cref{sec:modul} är hämtat ur dess
tabell~2.

Skriv sedan ner hur ni gick till väga, så att någon annan kan följa
spåret.  I det här materialet står den anteckningen som en kommentar
ovanför varje post i litteraturlistans källfil, med fasta fält:
\begin{description}
  \item[\texttt{CLAIM}] vilket påstående källan ska styrka;
  \item[\texttt{FOUND-VIA}] hur den hittades --- sökfråga och databas,
    eller \enquote{känt dokument};
  \item[\texttt{PICKED}] varför just den, framför andra träffar;
  \item[\texttt{QUOTE}] det ordagranna citatet, med sida;
  \item[\texttt{VERIFIED}] att fulltexten lästs, och var;
  \item[\texttt{COUNTER}] motsökningen och vad den gav;
  \item[\texttt{DATE}] när.
\end{description}
Det tar ett par minuter per källa, och det är det enda som gör det
möjligt att om ett år svara på frågan \enquote{hur vet du det?}

\section{Dra slutsatsen --- med förbehåll}

Svaret skrivs ut tillsammans med det som talar emot.  Förbehållen är
inte svagheter i svaret; de \emph{är} svaret, och de avgör hur
påståendet får användas i materialet.  Två exempel ur det här
materialet.  Beläggen för KISS visar att Kelly Johnson \emph{använde}
uttrycket, inte att han myntade det, och därför står det
\enquote{tillskrivs} i \cref{sec:kiss}.  Och Fleury studerade
procedurer i Pascal, inte moduler i Python; därför säger \cref{sec:modul}
att kontrasten mellan \mintinline{python}{typed_input.input_int} och
ett naket \mintinline{python}{input_int} är byggd \emph{mot samma slags}
regler som hon fann, inte att just den regeln är belagd.

Ett påstående kan också få strykas.  Det ligger nära till hands att
skriva att nybörjare \emph{ofta} döper sina filer till samma namn som en
modul (\cref{sec:skugga}).  Kursens sammanställning över dokumenterade
missuppfattningar täcker inte moduler, och ingen litteratursökning har
gjorts för den här föreläsningen, så påståendet står inte i materialet.
Det som står är vad körningen visar: att Python sedan version 3.13 ger
ett råd om att byta filnamn i själva felmeddelandet.  Vill någon
kontrollera hur vanligt misstaget är, är det en sökning som återstår att
göra.

\begin{table}[htbp]
  \centering
  \caption{Materialets påståenden, och var vart och ett är
    kontrollerat.}
  \label{tab:paastaenden}
  \footnotesize
  \begin{tabular}{@{}p{0.44\linewidth}p{0.48\linewidth}@{}}
    \toprule
    Påstående & Kontrollerat \\
    \midrule
    Att inte upprepa sig i ett program är en bra princip, och den
      härrör från Hunt och Thomas (DRY)
      & Redan belagt i \emph{Algoritmiskt tänkande}, bilaga E; posten
        är kopierad därifrån med sin proveniensanteckning \\
    SRP härrör från Robert C. Martin, och är en etablerad praktik
      & Redan belagt i \emph{Funktioner}, bilaga B \\
    KISS tillskrivs Kelly Johnson, och enkelhet är ett etablerat
      designvärde
      & Redan belagt i \emph{Funktioner}, bilaga C \\
    Vad en modul är, vad \mintinline{python}{import} gör, vad
      \mintinline{python}{__name__} innehåller, var Python letar efter
      moduler, vad \texttt{\_\_pycache\_\_} är, vad ett paket är, vad
      \texttt{pip} och virtuella miljöer gör, och vad PyPI är
      & Primärkällor: Pythons egen handledning och biblioteksreferens
        samt PyPI:s egen sida, hämtade och lästa 2026-09-03; beteendet
        bekräftat genom att bygget kör exempelprogrammen under Python
        3.13 \\
    Studenter bygger egna regler för var ett namn letas upp
      & Återanvänt från forskningsartikeln om missuppfattningar via
        \emph{Funktioner}; fulltexten läst på nytt för det här
        materialet \\
    Kontrast måste upplevas samtidigt, och en uppgift som pekar ut
      den kritiska aspekten lämnar inget att urskilja
      & Andrahandskälla: Martons bok, citerad som i kursens övriga
        material \\
    Nybörjare döper ofta filer till samma namn som en modul
      & \emph{Struket}: inte belagt, ingen sökning gjord; materialet
        påstår det inte \\
    \bottomrule
  \end{tabular}
\end{table}

\section{En anmärkning om språkmodeller}

I de föreläsningar där en litteratursökning gjordes gav sökningarna
hundratals träffar, och en språkmodell användes för att \emph{sortera}
dem --- hör träffen till ämnet över huvud taget, och åt vilket håll
pekar den? --- så att gallringen kunde granskas.  Lägg märke till
arbetsfördelningen: modellen sorterade, men varje källa som citeras är
läst och vald av en människa.  I det här materialet behövdes ingen
sortering, eftersom ingen ny sökning gjordes; men det material ni läser
är delvis skrivet med hjälp av en språkmodell, och samma regel gäller
då: modellen får hitta, sortera och formulera, men varje källa och varje
citat är kontrollerat mot originalet, och det är människan som står till
svars för innehållet.  Använd gärna språkmodeller för att hitta och
sortera; låt dem aldrig vara det sista ledet.
